% Generated draft: requires mathematical and source review.
\documentclass[11pt]{article}
\usepackage[margin=25mm]{geometry}
\usepackage{fontspec}
\setmainfont{Latin Modern Roman}
\usepackage{amsmath,amssymb,mathtools,unicode-math}
\setmathfont{Latin Modern Math}
\usepackage{enumitem,needspace,xurl,hyperref}
\setlength{\parindent}{0pt}
\setlength{\parskip}{5pt}
\setlength{\emergencystretch}{3em}
\newcommand{\N}{\mathbb{N}}
\newcommand{\Z}{\mathbb{Z}}
\newcommand{\Q}{\mathbb{Q}}
\newcommand{\R}{\mathbb{R}}
\newcommand{\C}{\mathbb{C}}
\newcommand{\F}{\mathbb{F}}
\newcommand{\A}{\mathbb{A}}
\newcommand{\PP}{\mathbb{P}}
\newcommand{\E}{\mathbb{E}}
\newcommand{\eps}{\varepsilon}
\newcommand{\dd}{\,\mathrm{d}}
\begin{document}
\section*{UnsolvedMath ID: 30003960}
\textbf{Source problem code:} catalogue entry
\subsection{Definitions and mathematical statement}
Let \(d\) be a positive integer and set \(X=\Z^d\), equipped with the metric \(\|x-y\|_1=\sum_{j=1}^d|x_j-y_j|\). The wobbling group is
\[
W(X)=\left\{g:X\to X:\ g\text{ is bijective and }\sup_{x\in X}\|g(x)-x\|_1<\infty\right\},
\]
with composition as its group operation and its natural action on \(X\). The positive-integer range for \(d\) makes explicit the usual lattice-dimensional interpretation; the supplied question does not explicitly specify that range.

Write \(\mathcal P_f(X)\) for the set of finite subsets of \(X\), including the empty set. The action \(G\curvearrowright X\) is called extensively amenable if there exists a finitely additive probability measure \(m\) on all subsets of \(\mathcal P_f(X)\) such that
\[
m(g\mathcal B)=m(\mathcal B)
\quad\text{for every }g\in G\text{ and }\mathcal B\subseteq\mathcal P_f(X),
\]
where \(g\mathcal B=\{gA:A\in\mathcal B\}\) and \(gA=\{g(x):x\in A\}\), and such that
\[
m\bigl(\{A\in\mathcal P_f(X):F\subseteq A\}\bigr)=1
\quad\text{for every finite }F\subseteq X.
\]

For the probabilistic formulation, let \(\mu\) be any finitely supported probability measure on \(W(X)\) satisfying \(\mu(g)=\mu(g^{-1})\) for every \(g\in W(X)\). Let \(g_1,g_2,\ldots\) be independent random elements with law \(\mu\), and use the explicit convention
\[
S_0=\operatorname{id}_X,\qquad S_k=g_k\circ\cdots\circ g_1,\qquad
O_n=\{S_k^{-1}(0):0\leq k\leq n\}.
\]
Thus \(O_n\) is the inverted orbit based at the origin under this convention. Neither the base point nor the product and endpoint conventions are specified in the supplied record; this definition supplies a conventional interpretation requiring comparison with the cited source. No assumption of laziness or of generation of the full wobbling group is imposed on \(\mu\).

\textbf{Question.} For each positive integer \(d\), is the natural action \(W(\Z^d)\curvearrowright\Z^d\) extensively amenable? The supplied record presents the following as an equivalent formulation: for every symmetric finitely supported probability measure \(\mu\) on \(W(\Z^d)\) and every real \(\varepsilon>0\), does there exist an integer \(N=N(d,\mu,\varepsilon)\geq1\) such that
\[
\forall n\in\Z,\quad n\geq N\ \Longrightarrow\
\mathbb P_\mu\bigl(|O_n|<\varepsilon n\bigr)\geq\exp(-\varepsilon n)?
\]
The equivalence is reported from the supplied record and is not independently verified here, including its compatibility with the explicit orbit convention above. The supplied dated assessment says that dimensions \(1\) and \(2\) are settled and that the case \(d\geq3\) and the all-measures formulation remain unresolved. Its phrase ``all-measures formulation'' does not specify an additional measure class beyond the symmetric finitely supported measures in the question; no broader class is inferred.
\subsection{Short English statement}
The question asks whether every wobbling action on \(\Z^d\) admits an invariant finitely additive probability measure on finite subsets that almost surely contain each prescribed finite set. The source also asks whether every symmetric finitely supported random walk has inverted orbits of size less than \(\varepsilon n\) with probability at least \(e^{-\varepsilon n}\) for all sufficiently large \(n\). Its dated curation identifies \(d\geq3\) as unresolved; this status and the claimed equivalence require review.
\subsection{Sources}
\begin{itemize}
\item Supplied local UnsolvedMath registry record, permanent ID \(30003960\), title: \emph{Extensive Amenability of Wobbling-Group Actions}; collection: Oberwolfach Reports Open Problems; exact source problem code: \texttt{catalogue entry}. No individual proposer is identified. The proposed year is recorded as 2018.
\item Supplied source citation: \emph{Scaling Limits in Models of Statistical Mechanics (2019)}, Oberwolfach Reports, pp.~2535--2581, DOI \texttt{10.4171/owr/2018/41}. Source and literature-source link, labelled ``Oberwolfach Report 41, extensive amenability'': \url{https://doi.org/10.4171/owr/2018/41}. The supplied background describes these reports as records of mathematical workshops and their open-problem discussions.
\item Supplied literature-source attribution: Garban, ``inverted orbits and extensive amenability'': \url{https://arxiv.org/abs/1804.01981}. This wording is the supplied reference label; complete bibliographic metadata and a pinpoint citation are not supplied.
\item Literature assessment checked at \texttt{2026-08-21}: the supplied curation labels the problem open and states that extensive amenability is settled in dimensions \(1\) and \(2\), while the \(d\geq3\) case and the all-measures formulation remain unresolved. This is supplied dated curation derived from the OWR record, not an independently verified literature assessment. The record itself requests independent verification of time-sensitive status and cited claims.
\item The registry labels the statement \texttt{corrected\_verified} and says that its mathematical content was preserved while notation, typography, and omitted definitions were repaired. These are supplied editorial assertions, not mathematical verification performed for this draft.
\end{itemize}
\end{document}
